% Propietario conservado: /Users/ruben/Documents/New project/output/EXTREMA_MEDIA_RAZON_RECIPROCIDAD_ES_EN_20260918/ARTICULO/ES/source/sections/01_hilbert_nonadico.tex
% Procedencia: integral 2249, cap. 23, pp. 601--608; propietario base_c21_sin_encabezado.tex.
% Recuperación y desarrollo de pruebas. Requiere el núcleo común APP--TRIT--TPK.
\section{Hilbert-space realization of the nonadic phase and preservation of the unit}
\label{x_reciprocidad:iv:hilbert}

The construction of Triadic Modular Holography first provides
the states, their transitions, and the restriction maps. The Hilbert-space
realization adds a linear structure and an inner
product to their discrete coordinates. This order determines what each operator represents:
the basis labels arise from the nonadic construction;
linear operations subsequently act on those labels. Genealogical
memory remains in the TPK state and accompanies the projections
used to represent it.

\subsection{Phase fiber, labeled basis, and inner product}

For an integer depth \(d\geq1\), let
\begin{equation}
 \mathcal C_{9,d}=(\ZZ/9\ZZ)^d,
 \qquad \mathcal H_d=\ell^2(\mathcal C_{9,d}),
 \qquad \mathcal A_d=B(\mathcal H_d).
 \label{x_reciprocidad:iv:eq:hilbert-finito}
\end{equation}
The inner product is
\(\langle f,g\rangle=\sum_x\overline{f(x)}g(x)\). The functions
\(\delta_x\) form an orthonormal basis labeled by the phase
coordinates. Thus \(\dim\mathcal H_d=9^d\) and
\(\mathcal A_d\cong M_{9^d}(\CC)\). The matrix algebra represents
this finite fiber; routes, quotients, and memory omitted by a
phase projection still belong to the starting state.

Coordinate factorization induces the unitary map
\begin{equation}
 \mathcal H_{d+1}\longrightarrow\mathcal H_d\otimes\CC^9,
 \qquad \delta_{(x,r)}\longmapsto\delta_x\otimes\delta_r.
 \label{x_reciprocidad:iv:eq:factorizacion-hilbert}
\end{equation}
Define the normalized trace
\(\tau_d(a)=9^{-d}\operatorname{Tr}(a)\). Finite counting and
normalization are distinguished: \(\operatorname{Tr}(I)=9^d\), whereas
\(\tau_d(I)=1\). Likewise, the group \((\ZZ/9\ZZ)^3\) and the
vector space \(\FF_3^6\) both have \(729\) elements, but
retain different algebraic laws. Realization of a phase
fiber does not replace the ternary encoding of APP emissions.

\subsection{Translation and phase: finite Weyl representation}

In the complex realization subsequent to the constructions of
\(\pi\) and the structure \(\mathrm i^2=-1\), fix
\(\omega=\exp(2\pi\mathrm i/9)\). On \(\mathcal H_1\), let
\begin{equation}
 (Xf)(r)=f(r-1),\qquad (Zf)(r)=\omega^r f(r).
 \label{x_reciprocidad:iv:eq:weyl-operadores}
\end{equation}
The first operation shifts the phase coordinate; the second
evaluates it through a character. Both are unitary and satisfy
\begin{equation}
 X^9=Z^9=I,\qquad ZX=\omega XZ.
 \label{x_reciprocidad:iv:eq:weyl-relacion}
\end{equation}
Indeed, \((ZXf)(r)=\omega^r f(r-1)\), whereas
\((XZf)(r)=\omega^{r-1}f(r-1)\). The character appears here as a
realization of the phase already constructed and does not select the TPK
numerical registers.

% RESULTADO_RECUPERADO: 14_numero_heisenberg_y_d108.tex:741--788.
% Prueba algebraica desarrollada con la convención X^a Z^b ya fijada aquí.
\subsubsection{Oriented area and central extension of translations}
\label{x_reciprocidad:iv:weyl-area-extension}

The translation support is \(X_9=(\ZZ/9\ZZ)^2\).
The mixed APP curvature constructed in
Section~\ref{x_reciprocidad:iv:extension-curvatura-mixta} assigns \(+1\) to the
oriented plaquette of the sum--product pair and \(-1\) to the
pair in the opposite order. Its bilinear extension to two displacements
\(u=(a,b)\) and \(v=(c,d)\) is the alternating form
\begin{equation}
 \sigma_{\mathrm{ar}}(u,v)=ad-bc\quad\text{in }\ZZ/9\ZZ.
 \label{x_reciprocidad:iv:eq:weyl-area-orientada}
\end{equation}
Indeed, alternation makes the evaluation on two displacements
along the same axis vanish. The values on the oriented basis are
\(\sigma_{\mathrm{ar}}((1,0),(0,1))=1\) and
\(\sigma_{\mathrm{ar}}((0,1),(1,0))=-1\); distributing both
arguments gives \eqref{x_reciprocidad:iv:eq:weyl-area-orientada}. The same formula is
the sum of plaquettes of any oriented cellular chain realizing
that parallelogram: the interior edges cancel. Changing an
integer representative by a multiple of nine preserves its publication
in \(\ZZ/9\ZZ\).

The order \(X^aZ^b\) of the already defined operators fixes a section
of the central extension. Denote its composition factor by
\begin{equation}
 c_{\mathrm W}(u,v)=bc\quad\text{in }\ZZ/9\ZZ.
 \label{x_reciprocidad:iv:eq:weyl-cociclo-polarizado}
\end{equation}
This cocycle preserves the order of the two operations, while
\(\sigma_{\mathrm{ar}}\) preserves the orientation of the parallelogram.

\begin{proposition}[Central extension and commutator sign]
\label{x_reciprocidad:iv:prop:heisenberg-orientado}
The set \(\mathfrak H_9=(\ZZ/9\ZZ)^3\), with product
\begin{equation}
 (a,b,t)(c,d,s)=(a+c,b+d,t+s+bc),
 \label{x_reciprocidad:iv:eq:heisenberg-ley}
\end{equation}
is a group of order \(729\). The map
\[
 \varrho(a,b,t)=\omega^tX^aZ^b
\]
is a faithful unitary representation on \(\mathcal H_1\).
With the convention \([g,h]=ghg^{-1}h^{-1}\), one has
\begin{equation}
 [\varrho(a,b,0),\varrho(c,d,0)]
 =\omega^{bc-ad}I
 =\omega^{-\sigma_{\mathrm{ar}}((a,b),(c,d))}I.
 \label{x_reciprocidad:iv:eq:weyl-conmutador-orientado}
\end{equation}
The center and the translation group form the exact sequence
\[
 0\longrightarrow\ZZ/9\ZZ
 \xrightarrow{\ t\mapsto(0,0,t)\ }\mathfrak H_9
 \xrightarrow{\ (a,b,t)\mapsto(a,b)\ }X_9
 \longrightarrow0.
\]
\end{proposition}
\begin{proof}
For \(w=(e,f)\), the cocycle identity reads
\[
 c_{\mathrm W}(u,v)+c_{\mathrm W}(u+v,w)
 =bc+(b+d)e
 =de+b(c+e)
 =c_{\mathrm W}(v,w)+c_{\mathrm W}(u,v+w).
\]
It proves associativity of \eqref{x_reciprocidad:iv:eq:heisenberg-ley}; the identity
is \((0,0,0)\), and the inverse of \((a,b,t)\) is
\((-a,-b,-t+ab)\). The three residues are free, giving order
\(9^3=729\). From \(ZX=\omega XZ\) it follows that
\[
 (X^aZ^b)(X^cZ^d)=\omega^{bc}X^{a+c}Z^{b+d},
\]
which proves the representation property. If
\(\omega^tX^aZ^b=I\), its action on \(\delta_r\) first forces
\(a=0\). Evaluation at \(r=0\) gives \(t=0\), and
evaluation at \(r=1\) gives \(b=0\), since \(\omega\) has order
nine. The representation is therefore faithful. Comparing the
products in both orders gives the factor \(bc-ad\).
A central element must satisfy \(bc-ad=0\) for every \((c,d)\).
Taking successively \((c,d)=(1,0),(0,1)\) yields
\(b=a=0\). This identifies the center and proves exactness.
\end{proof}

The Wilson phase for the orientation
\eqref{x_reciprocidad:iv:eq:weyl-area-orientada} is
\(\omega^{\sigma_{\mathrm{ar}}(u,v)}\); the commutator of the
ordered section \(X^aZ^b\) publishes its inverse. In particular,
\(c_{\mathrm W}(u,v)-c_{\mathrm W}(v,u)
=-\sigma_{\mathrm{ar}}(u,v)\).
This identity relates two specified conventions, without identifying
the polarized cocycle with the alternating form. The extension belongs
to the translation group of the APP support. Its representation acts on
the phase fibre; the path register and carries retain
their residence in the TPK state that produced that fibre. The order \(729\)
of the group and the cardinality \(729\) of the ternary ambient space retain the
distinct types established in \eqref{x_reciprocidad:iv:eq:hilbert-finito}.
The centre \(\ZZ/9\ZZ\) of this extension records the phase
of the commutator. The memory \(C_{12}\) of the sequence
\(0\to C_{12}\to C_{108}\to C_9\to0\), constructed in
proposition~\ref{x_reciprocidad:prop:k-extension}, belongs to the dodecaphase
calendar. Each extension retains its own composition law and
projection.

\begin{proposition}[Algebra generated by phase and translation]
\label{x_reciprocidad:iv:prop:weyl-completo}
The \(81\) operators \(X^aZ^b\), with \(0\leq a,b<9\), form
a basis of \(\mathcal A_1=M_9(\CC)\). In particular, the
representation \eqref{x_reciprocidad:iv:eq:weyl-operadores} is irreducible.
\end{proposition}
\begin{proof}
If \(a\neq0\pmod9\), the matrix \(X^aZ^b\) has no diagonal
entries. If \(a=0\), its trace is
\(\sum_{r=0}^8\omega^{br}\), equal to \(9\) for \(b=0\) and
\(0\) otherwise. By \eqref{x_reciprocidad:iv:eq:weyl-relacion}, the
product \((X^aZ^b)^*X^{a'}Z^{b'}\) is a factor of modulus one
times \(X^{a'-a}Z^{b'-b}\). Therefore,
\[
 \tau_1\bigl((X^aZ^b)^*X^{a'}Z^{b'}\bigr)
 =\delta_{a,a'}\delta_{b,b'}.
\]
These are \(81\) linearly independent matrices in a space of
dimension \(81\). They generate the whole matrix algebra, whose commutant
consists of scalars; irreducibility follows.
\end{proof}

\begin{theorem}[Uniqueness with fixed primitive central character]
\label{x_reciprocidad:iv:thm:weyl-unicidad}
Every irreducible unitary representation of the group presented by
\(X^9=Z^9=c^9=I\), \(c\) central, and \(ZX=cXZ\), in which
\(c\) acts as \(\omega I\), is unitarily equivalent to
\eqref{x_reciprocidad:iv:eq:weyl-operadores}.
\end{theorem}
\begin{proof}
Let \(v\neq0\) be an eigenvector of \(Z\), with eigenvalue \(\lambda\).
Since \(Z^9=I\), \(\lambda\) is a power of \(\omega\).
The vectors \(X^kv\), \(0\leq k<9\), have eigenvalues
\(\omega^k\lambda\), distinct because the character is primitive.
They are orthogonal, and their span is invariant under \(X\),
\(Z\), and the center. Irreducibility forces it to be the entire
space. Reordering the basis to begin with eigenvalue \(1\),
\(Z\) is diagonal with entries \(1,\omega,\ldots,\omega^8\),
and \(X\) is cyclic translation. Normalizing the basis
produces the asserted unitary equivalence.
\end{proof}

Primitivity of the character and irreducibility are part of the
statement. At depth \(d\), operators acting on different
factors commute, and tensor products of the preceding bases
generate \(\mathcal A_d\). The phase structure is thus
represented at each level without identifying phase return with
return of the complete genealogical state.

\subsection{Two lifts and two different normalizations}

Factorization \eqref{x_reciprocidad:iv:eq:factorizacion-hilbert} permits preserving
the entire new fiber or fixing one of its labels. For
\(p_j=|\delta_j\rangle\langle\delta_j|\), define
\begin{equation}
 \iota_d(a)=a\otimes I_9,
 \qquad \jmath_{d,j}(a)=a\otimes p_j.
 \label{x_reciprocidad:iv:eq:elevaciones}
\end{equation}
\begin{proposition}[Preservation of unit and trace]
\label{x_reciprocidad:iv:prop:unidad-traza}
The inclusion \(\iota_d\) is unital and preserves normalized trace.
The inclusion \(\jmath_{d,j}\) is a corner embedding and
satisfies
\[
 \tau_{d+1}(\iota_d(a))=\tau_d(a),\quad \iota_d(I)=I,
 \qquad
 \tau_{d+1}(\jmath_{d,j}(a))=\frac{\tau_d(a)}9,
 \quad \jmath_{d,j}(I)=I\otimes p_j.
\]
\end{proposition}
\begin{proof}
The trace of a tensor product is the product of the traces.
Since \(\operatorname{Tr}(I_9)=9\) and
\(\operatorname{Tr}(p_j)=1\), dividing by \(9^{d+1}\)
produces both formulas. The images of the identity are read
directly from \eqref{x_reciprocidad:iv:eq:elevaciones}.
\end{proof}

The unital inclusion preserves the set of \(9\) phases and its
global normalization. The corner preserves a marked continuation and
its support projector. Both operations are legitimate, but
transmit different data to the next depth.

\subsection{Inductive closure and tracial realization}

\begin{theorem}[Nonadic uniformly hyperfinite algebra]
\label{x_reciprocidad:iv:thm:uhf}
The norm closure of the union through the inclusions \(\iota_d\)
is
\[
 \mathcal A_{9^\infty}
 =\varinjlim(\mathcal A_d,\iota_d)
 \cong\bigotimes_{n=1}^{\infty}M_9(\CC).
\]
It is a UHF algebra of type \(9^\infty=3^\infty\), with a unique
normalized trace \(\tau\), whose restriction to \(\mathcal A_d\)
is \(\tau_d\).
\end{theorem}
\begin{proof}
The inclusions are injective and unital. The closure of the increasing
union of matrix algebras is precisely the UHF construction.
The compatibility in Proposition \ref{x_reciprocidad:iv:prop:unidad-traza}
defines \(\tau\) on the union; positivity and norm one allow
its extension to the closure. Every normalized trace on the limit restricts
to the unique normalized trace on each matrix algebra. Since the union
is dense, those restrictions determine the trace uniquely.
\end{proof}

Let \(\mathcal H_\tau\) be the completion of
\(\mathcal A_{9^\infty}\) for
\(\langle a,b\rangle_\tau=\tau(a^*b)\), and let \(\pi_\tau\)
be its representation by left multiplication. The vector defined
by the identity is denoted \(\Omega_\tau\).

\begin{theorem}[Tracial factor of complete refinement]
\label{x_reciprocidad:iv:thm:factor-tracial}
The weak closure
\(\mathcal R_9=\pi_\tau(\mathcal A_{9^\infty})''\)
is a hyperfinite factor of type \(II_1\). The nonadic
inclusions preserve its unit and normalized trace.
\end{theorem}
\begin{proof}
The product trace extends to a faithful finite normal trace
in the tracial representation. To see that the center is trivial,
let \(E_d\) be the tracial expectation onto \(\pi_\tau(\mathcal A_d)\).
On a tensor of length \(n>d\), it takes the trace of the last
\(n-d\) factors. These maps extend continuously
to the corresponding orthogonal projections in \(L^2\).
If \(z\) is central, \(E_d(z)\) commutes with
\(\pi_\tau(\mathcal A_d)\), by bimodularity of \(E_d\).
It is therefore scalar and equal to \(\tau(z)I\). Density of the
matrix union implies \(E_d(z)\to z\) in \(L^2\), hence
\(z=\tau(z)I\). The algebra is thus a finite factor.
It contains unital matrix algebras of unbounded sizes \(9^d\),
so it is not a finite matrix factor. It is of type \(II_1\).
It is hyperfinite because the same union of matrices is weakly
dense. Preservation of unit and trace follows from each finite
inclusion and persists upon taking closures.
\end{proof}

\subsection{Integer ranks and continuous tracial dimension}

In \(\mathcal A_d\), a projection \(P\) has normalized
dimension \(\tau_d(P)=\operatorname{rank}(P)/9^d\).
Ranks remain integers at every finite depth. Their
compatible normalizations allow the following passage to the limit.

\begin{proposition}[Realization of each tracial dimension]
\label{x_reciprocidad:iv:prop:dimension-tracial}
For each \(t\in[0,1]\), there exists a projection
\(P_t\in\mathcal R_9\) with \(\tau(P_t)=t\).
\end{proposition}
\begin{proof}
Let \(m_d=\lfloor9^dt\rfloor\). Then
\(0\leq m_{d+1}-9m_d\leq8\). After choosing a diagonal projection
of rank \(m_d\), its inclusion has rank \(9m_d\);
add \(m_{d+1}-9m_d\) diagonal positions in its complement.
This yields an increasing sequence of projections within
\(\mathcal R_9\). Its strong limit is a projection \(P_t\),
and normality of the trace gives
\[
 \tau(P_t)=\lim_{d\to\infty}\frac{\lfloor9^dt\rfloor}{9^d}=t.
\]
\end{proof}

Here \(t\) denotes the dimension coordinate represented
in the closure, not an input to the APP generator. The assertion
describes the scope of the resulting tracial realization. To
compare it with a particular observable, it remains necessary to
identify the projection and the genealogical information it preserves.

By contrast, the isometries
\(V_{d,j}\xi=\xi\otimes\delta_j\) realize
\(\jmath_{d,j}(a)=V_{d,j}aV_{d,j}^*\). In the limit of a
marked route, the finite spaces are dense. The corners
approximate all finite-rank operators; their norm closure
is \(\mathcal K(\mathcal H_\infty)\), and their bicommutant is
\(B(\mathcal H_\infty)\), of type \(I_\infty\). This
difference preserves the operational content of the two choices
in \eqref{x_reciprocidad:iv:eq:elevaciones}: selection of one continuation and
unital prolongation of the entire fiber have different limits.

The Hilbert-space realization, its phase operators, and the tracial
closure are thus constructed in this section, with their inclusion
choices explicit. They constitute the
operator-theoretic support to be used subsequently for coordinate
interchange and the compact Hamiltonian.
